% ======================================================================
% SECTION 5: LIMITATIONS AND FUTURE WORK
% Two thematic blocks: per-simulation limits, and future-work direction.
% ======================================================================

\subsection{Per-Simulation Limitations}
\label{sec:limits-per-sim}

The four simulations have distinct limitations that reflect their
design tradeoffs.

\textbf{Simulation 1 (continuous RTCT, site).} The Continuity-PSL
metric introduced in \S\ref{sec:results-sim1} is approximate. The
rolling 24-hour weighted mean is tuned by hand rather than calibrated
on real patient outcomes, so the absolute C-PSL value is not directly
comparable to a TRIPOD+AI-compliant predictive index. Hours 56
through 83 contain approximated diagrams under \path{extra-hours/}
and are excluded from the full paper text per Methods
(\S\ref{sec:methods-inputs}). No separate code agents for local
compute were provided in this simulation; the entire output is text
only, which limits downstream automation.

\textbf{Simulation 2 (single-patient journey, site).} The simulation
covers a single patient (PAT-2026-0042) and is not as extensive as
the multi-patient continuous trial in Simulation 1. There is no
national network model and no multi-site aggregation. The single-
patient orientation limits population-level prediction calibration to
TRIPOD+AI and CREMLS standards \cite{Collins2024TRIPODAI,
ElEmam2024CREMLS}.

\textbf{Simulation 3 (24-hour autonomous sponsor).} The 24-hour
scope, while complete, does not exercise multi-day continuity. The
sponsor decisions were not reviewed by a real pharmaceutical sponsor
and the agent decision logic is synthetic (rule-based dispatch with
LLM-generated rationale text). The 5-endpoint FastAPI server is a
demonstration scaffold, not a production-grade sponsor information
system.

\textbf{Simulation 4 (168-hour 7-day extension, sponsor).} The 7-day
run was simulated at high speed, not at wall-clock 7 days. Escalations
and patient outcomes were generated by deterministic agent rules
rather than by clinical observation. The local verification on a 4 GB
laptop was partial - hour 134 paused for thermal throttling and
resumed only after intervention - and downstream analytics must
account for the partial-verification status.

\textbf{Cross-simulation limitations (all four).} Three concerns apply
to every simulation in this paper:

\begin{itemize}[leftmargin=*]
\item None of the four simulations was tested on real patients. All
patient identifiers, vital signs, and outcomes are synthetic or
illustrative. The work demonstrates a computational and software
capability, not a clinical capability.

\item Additional generations may vary in output type. Claude Code
Opus 4.7 Max output is non-deterministic; a re-run of any simulation
could produce different ASCII diagrams, different agent decisions, or
different PSL trajectories. The synthesis presented in Results
reflects one realization, not a population of realizations.

\item The four simulations primarily show the computational and
software benefits that LLM-scale repository context offers over
current oncology-trial AI. They do not replace TRIPOD+AI or CREMLS
reporting requirements for any deployment-grade extension
\cite{Collins2024TRIPODAI, ElEmam2024CREMLS}.
\end{itemize}

The single most important limitation across the four simulations is
the absence of real-patient validation. The synthetic patients,
synthetic robot logs, and synthetic agent decisions reflect carefully
constructed scenarios, but they do not yet provide evidence that the
LLM-scale repository context approach improves real patient outcomes.
Bridging that gap is the subject of the future work tracks below.

\subsection{Future Work}
\label{sec:limits-future}

Two distinct future-work tracks extend the computational advantages
demonstrated here into real-world applications with real patients.

\textbf{Track A: Single big model performing all tasks.} A future
generation of Claude Code (or an equivalent foundation-model agent)
holds the entire trial repository in context and emits all patient
predictions, sponsor decisions, and FDA-bound signals from a single
inference loop. This track is the natural extension of Simulation 4
- take the 168-hour autonomous sponsor and connect its output stream
to a real-time clinical workflow under TRIPOD+AI and CREMLS reporting
standards \cite{Collins2024TRIPODAI, ElEmam2024CREMLS}. Track A is
architecturally simple: one model, one inference loop, one
accountability surface for regulators. Its weaknesses are higher
inference cost per decision (a 1M-token context call is not cheap)
and a single point of regulatory accountability (when the model
makes an error, the entire decision trail must be reviewed at once).

\textbf{Track B: Single big model that creates smaller agents
performing smaller tasks locally.} A future generation uses the 1M-
token-context model only for orchestration and routing, while
delegating per-task prediction (mortality, toxicity, response) to
smaller specialized agents that can run on-site on hardware as
constrained as the Core i5-6200U used in Simulation 4. This track is
the natural extension of Simulation 1's site-level architecture and
Simulation 4's local verification. The small agents run on-site
(handling PHI without leaving the site), the big model coordinates
across sites and streams endpoints to the FDA in real time. Track B
is architecturally more complex and distributes regulatory
accountability across many small agents - each agent must satisfy
its own TRIPOD+AI or CREMLS reporting floor - but it lowers per-site
inference cost and improves the security posture for PHI handling.

\begin{table}[H]
\centering
\caption{Track A versus Track B trade-offs for a future RTCT-aligned
oncology trial AI deployment.}
\label{tab:future-tracks}
{\renewcommand{\arraystretch}{1.15}
\begin{tabular}{>{\raggedright\arraybackslash}p{3.5cm} >{\raggedright\arraybackslash}p{5cm} >{\raggedright\arraybackslash}p{6.6cm}}
\toprule
\textbf{Property} & \textbf{Track A: Big Single Model} & \textbf{Track B: Big Model + Small Agents} \\
\midrule
Architecture           & 1 model, 1 inference loop                & 1 orchestrator plus N specialist agents       \\
Inference cost         & Higher per decision                       & Lower per site, higher overall design cost     \\
Reg. accountability    & Single surface                            & Distributed; each agent has its own floor      \\
PHI handling           & Cloud transit required                    & Local agents handle PHI on-site                \\
Hardware floor         & Cloud (data-center)                       & Site (Core i5-6200U / 4 GB feasible)           \\
Update cycle           & Single model upgrade                      & Per-agent upgrade, with orchestrator swap      \\
Extension target       & Simulation 4 + clinical link              & Simulation 1 site network, per-task agents    \\
\bottomrule
\end{tabular}}
\end{table}

The choice between Track A and Track B is a research question, not a
settled matter. Both tracks should be prototyped before a real-
patient trial is designed, and the comparison itself becomes a
research output. The four simulations in this paper give each track
a starting point: Track A starts from Simulation 4, Track B starts
from Simulations 1 and 4 paired together.

\textbf{Concrete future deliverables.} Three concrete future
deliverables follow from the Track A / Track B comparison:

\begin{enumerate}[leftmargin=*]
\item A TRIPOD+AI / CREMLS-compliant retrospective validation of one
of the four simulations on a real-patient cohort. The most natural
candidate is Simulation 1's hour-resolution patient narrative, which
could be retrospectively scored against a real EHR-linked oncology
cohort to test whether the LLM-scale narrative carries useful
prognostic signal beyond the structured EHR features.

\item A public benchmark comparing the LLM-scale repository context
approach to the supervised baselines named in Background-A and
Background-B (Manz 2020, SHIELD-RT, SCORPIO, PROGPATH, AIM-LCpro).
The benchmark would fix a prediction task (180-day mortality, AE
prediction, or ICI response), supply both structured features and
the unstructured repository context, and report AUROC and C-index
under TRIPOD+AI reporting standards.

\item An FDA-aligned RTCT pilot submission that uses the artifacts
from Simulation 4 as inputs. The submission would propose pairing
Simulation 1 (site) with Simulation 4 (sponsor) as the technical
framework, demonstrate the local-verification proof on a Core i5-
6200U laptop, and submit the GitHub commit log as the audit trail.
\end{enumerate}

\textbf{Local AI compute for safety and security.} The future-work
direction with the most direct impact on patient safety and security
is a future Claude Code or competing AI local instance that
combines the power of cloud compute (1M token context) with the
flexibility and security of local computing. A local instance would
let an oncology trial site keep PHI on-site, run specialist agents
on local hardware (Core i5-6200U or similar), and connect to the
cloud orchestrator only for repository-scale reasoning. This hybrid
architecture is the practical realization of Track B and is the
direction in which the four simulations here naturally point.

The hybrid local/cloud architecture also addresses a regulatory
concern that the cloud-only approach does not. A trial sponsor that
processes PHI on local hardware can produce an audit trail that an
FDA inspector can verify physically (a local hash plus a local
commit), without requiring transparency cooperation from a cloud
provider. The Simulation 4 local-verification run on a Core i5-6200U
demonstrates that this trail is feasible on commodity hardware. A
future RTCT pilot that uses local compute for PHI and cloud compute
only for non-PHI orchestration would inherit this audit advantage.

\textbf{Reporting standards integration.} Any future deliverable must
adopt the TRIPOD+AI and CREMLS reporting frameworks
\cite{Collins2024TRIPODAI, ElEmam2024CREMLS} as the floor. Both
frameworks anticipate the inclusion of LLM-scale prediction systems,
but neither yet fully specifies the reporting elements for a
repository-scale prediction system that emits multiple artifact types
per inference pass. A future revision of TRIPOD+AI or CREMLS that
adds reporting elements for repository-scale LLM prediction would
make the four-simulation set directly comparable to the supervised
baselines under a shared reporting floor.
